% Reconstructed text from V5 PDF page 39; non-native source
\par\smallskip\noindent\rule{\textwidth}{0.4pt}\par{\sffamily\footnotesize Original V5 page 39 -- recovered text}\par\begin{flushleft}\footnotesize\ttfamily
\par\smallskip
\par\smallskip
\noindent Diagnostics and Serving Infrastructure for MAS 39\par
\par\smallskip
\par\smallskip
\par\smallskip
\noindent depend on, coverage may degrade silently. The risk is acute for ShieldAgent [8], which main-\par
\noindent tains short-term KV-cache memory “discarded after verification” — if the serving system evicts\par
\noindent this before verification completes, the safety check is bypassed. Characterizing this coupling\par
\noindent requires controlled experiments varying memory pressure while measuring safety coverage met-\par
\noindent rics (false negative rate, bypass frequency).\par
\par\smallskip
\noindent 7.2 Diagnostic Overhead Characterization\par
\noindent A prerequisite question for the validation protocol is whether diagnostic tracing itself introduces\par
\noindent meaningful overhead. We benchmarked four tracing modes in a simulated multi-agent pipeline:\par
\noindent no tracing, string-formatted logging, structured OTel-style spans, and causal DAG construction.\par
\noindent  All tracing modes add estimated less than 3\% CPU overhead relative to the untraced baseline\par
\noindent in our simulated multi-agent pipeline (single-threaded Python, 3 agents, 50 conversations; aver-\par
\noindent aged over 3 runs). Structured OTel-style spans are slightly more efficient than string-formatted\par
\noindent logging (the formatting overhead of string concatenation exceeds the metadata overhead of\par
\noindent span construction). Causal DAG mode uses approximately 2× the memory of structured spans\par
\noindent at high verbosity, because it maintains the full parent-child relationship graph in memory rather\par
\noindent than streaming events to a sink. These measurements reflect CPU-side tracing cost in isolation;\par
\noindent GPU-level overhead from shared memory contention with KV cache is not captured by this\par
\noindent benchmark.\par
\noindent  This result establishes that tracing overhead in isolation is negligible on CPU. The diagnostic-\par
\noindent infrastructure interaction gap is therefore not about the raw cost of trace capture — it is about\par
\noindent resource contention with serving at the GPU level, where KV cache memory and diagnostic\par
\noindent evidence may compete for the same limited pool. No existing benchmark measures this GPU-\par
\noindent level contention, and the overhead benchmark confirms that the right measurement target is\par
\noindent GPU memory allocation under concurrent diagnostic and serving workloads, not CPU tracing\par
\noindent cost.\par
\par\smallskip
\noindent 7.3 Validation Protocol for the Shared-Budget Hypothesis\par
\noindent We propose a concrete experimental protocol for testing whether integrating diagnostic evi-\par
\noindent dence retention into the serving-side eviction objective improves end-to-end failure attribution\par
\noindent without unacceptable throughput loss.\par
\par\smallskip
\noindent Setup. Deploy a multi-agent system with: (a) a KV-cache serving backend with configurable\par
\noindent eviction policy (SGLang or Mooncake as the base); (b) a diagnostic trace store with configurable\par
\noindent retention budget; (c) injected failures at known locations for ground-truth attribution.\par
\par\smallskip
\noindent Independent variables.\par
\noindent  1. Retention budget (b): the fraction of available memory allocated to diagnostic evidence\par
\noindent  retention, varied from 0\% (pure serving) to 50\% (equal split).\par
\noindent  2. Eviction policy: baseline (LRU, serving-only objective) vs. diagnostic-aware (serving +\par
\noindent  diagnostic-value term).\par
\noindent  3. Workload intensity: low (single-agent, sequential tasks) to high (multi-agent, concurrent\par
\noindent  workflows with prefix sharing).\par
\end{flushleft}
